\documentclass[12pt,letterpaper]{article}

\usepackage[utf8]{inputenc}
\usepackage[T1]{fontenc}
\usepackage[spanish]{babel}
\usepackage{geometry}
\usepackage{setspace}
\usepackage{array}
\usepackage{longtable}
\usepackage{booktabs}
\usepackage{hyperref}
\usepackage{xcolor}
\usepackage{listings}
\usepackage{fancyhdr}
\usepackage{titlesec}
\usepackage{enumitem}

\geometry{margin=2.5cm}
\onehalfspacing
\setlength{\parindent}{1.25cm}
\setlength{\parskip}{0.25em}
\emergencystretch=3em
\sloppy
\hypersetup{
    colorlinks=true,
    linkcolor=blue,
    urlcolor=blue,
    citecolor=blue
}

\pagestyle{fancy}
\fancyhf{}
\lhead{SentinelMail AI}
\rhead{EIF420O - Inteligencia Artificial}
\cfoot{\thepage}

\titleformat{\section}{\normalfont\Large\bfseries}{\thesection.}{0.5em}{}
\titleformat{\subsection}{\normalfont\large\bfseries}{\thesubsection.}{0.5em}{}
\titleformat{\subsubsection}{\normalfont\normalsize\bfseries}{\thesubsubsection.}{0.5em}{}

\lstdefinestyle{pythonstyle}{
    language=Python,
    basicstyle=\ttfamily\small,
    keywordstyle=\color{blue},
    stringstyle=\color{teal},
    commentstyle=\color{gray},
    showstringspaces=false,
    breaklines=true,
    frame=single,
    columns=fullflexible
}

\begin{document}

\begin{titlepage}
    \centering
    \vspace*{1cm}
    {\Large Universidad Nacional\par}
    {\large Sede Regional Brunca -- Campus Coto\par}
    {\large Bachillerato en Ingeniería de Sistemas de Información\par}
    \vspace{1.5cm}
    {\large EIF420O -- Inteligencia Artificial\par}
    \vspace{2cm}
    {\Huge\bfseries SentinelMail AI\par}
    \vspace{0.4cm}
    {\Large Sistema inteligente de detección de phishing en correos electrónicos y URLs mediante red neuronal, Naive Bayes, árbol de decisión y sistema experto\par}
    \vspace{2cm}
    {\large Proyecto final\par}
    \vspace{1cm}
    \begin{tabular}{rl}
        Estudiantes: & Marco Ugalde Cortes \\
                     & Jhashua Canton Ruiz \\
        Profesor: & Ing. Jeremy Elizondo Castro \\
        Curso: & Inteligencia Artificial \\
        Ciclo lectivo: & I Ciclo 2026 \\
        Fecha: & 12 de junio de 2026 \\
    \end{tabular}
    \vfill
    {\large Costa Rica\par}
\end{titlepage}

\tableofcontents
\newpage

\section{Introducción}

El phishing es una de las amenazas más frecuentes en el uso cotidiano del correo electrónico. Consiste en engañar a una persona para que entregue credenciales, datos bancarios, información personal o acceso a sistemas por medio de mensajes que aparentan provenir de instituciones legítimas. Aunque muchas campañas de phishing son técnicamente simples, siguen siendo efectivas porque combinan elementos de ingeniería social, apariencia institucional, enlaces engañosos, urgencia psicológica y, en algunos casos, archivos adjuntos peligrosos.

El presente proyecto, denominado \textbf{SentinelMail AI}, propone una aplicación funcional para analizar correos electrónicos y URLs con el objetivo de estimar su riesgo de phishing. El sistema permite ingresar manualmente información del correo, analizar URLs, revisar cabeceras relevantes, cargar archivos con formato \texttt{.eml}, identificar señales en adjuntos y generar un resultado comprensible para el usuario. La salida principal del sistema es una clasificación de riesgo: \textbf{BAJO}, \textbf{MEDIO} o \textbf{ALTO}, acompañada por un porcentaje de riesgo y una explicación de las señales detectadas.

La motivación principal del proyecto es demostrar que las técnicas de inteligencia artificial vistas en el curso pueden aplicarse a un problema práctico y actual. La detección de phishing es adecuada para este propósito porque combina datos textuales, datos numéricos, reglas de conocimiento experto y patrones que pueden aprenderse mediante clasificación supervisada. Por esta razón, el proyecto no utiliza una única técnica, sino cuatro enfoques complementarios: una red neuronal artificial, un clasificador Naive Bayes, un árbol de decisión y un sistema experto basado en reglas de ciberseguridad.

El sistema fue desarrollado en Python y se implementó con bibliotecas estándar del lenguaje. Los algoritmos centrales de inteligencia artificial fueron programados manualmente por el grupo, sin depender de bibliotecas de aprendizaje automático como Scikit-learn, TensorFlow o PyTorch. Esta decisión responde directamente al requisito técnico del proyecto: comprender los fundamentos teóricos de las técnicas y llevarlos a una implementación propia.

\section{Descripción del problema}

El problema abordado consiste en clasificar correos electrónicos o enlaces como legítimos, sospechosos o probablemente maliciosos. Esta clasificación no es trivial porque un correo de phishing puede parecer normal a simple vista. Un atacante puede copiar el estilo de una institución, usar dominios visualmente parecidos, ocultar enlaces detrás de botones HTML, incluir archivos comprimidos o aprovechar cabeceras de correo inconsistentes.

Desde el punto de vista computacional, el problema puede formularse como una tarea de clasificación binaria y de puntuación de riesgo. La clase positiva corresponde a phishing y la clase negativa corresponde a correo legítimo. Sin embargo, para que la aplicación sea útil en una demostración, el resultado no se presenta únicamente como ``phishing'' o ``legítimo'', sino como un nivel de riesgo graduado. Esto permite explicar casos intermedios en los que existen algunas señales sospechosas, pero no suficientes para afirmar con alta confianza que el mensaje es malicioso.

El análisis se basa en distintas fuentes de información:

\begin{itemize}
    \item El asunto del correo.
    \item El cuerpo del mensaje.
    \item La URL principal.
    \item El remitente visible.
    \item Las cabeceras \texttt{Reply-To}, \texttt{Return-Path} y \texttt{Authentication-Results}.
    \item Los enlaces extraídos de contenido HTML o texto plano.
    \item Los archivos adjuntos incluidos en mensajes \texttt{.eml}.
\end{itemize}

La importancia de usar varias fuentes se debe a que el phishing no siempre se manifiesta en una sola señal. Un correo puede usar HTTPS y aun así ser malicioso; otro puede no contener adjuntos, pero solicitar credenciales; otro puede tener un texto aparentemente normal, pero fallar autenticaciones SPF, DKIM o DMARC. Por ello, SentinelMail AI combina modelos aprendidos y reglas expertas para obtener una evaluación más completa.

\section{Objetivos}

\subsection{Objetivo general}

Desarrollar una aplicación funcional capaz de analizar correos electrónicos y URLs para clasificar el riesgo de phishing mediante técnicas de inteligencia artificial implementadas manualmente.

\subsection{Objetivos específicos}

\begin{itemize}
    \item Extraer características relevantes de URLs, mensajes, cabeceras y adjuntos.
    \item Implementar una red neuronal artificial para clasificación binaria.
    \item Implementar un clasificador Naive Bayes para análisis textual.
    \item Implementar un árbol de decisión basado en impureza Gini.
    \item Implementar un sistema experto con reglas de ciberseguridad.
    \item Combinar los resultados de las técnicas mediante una política de puntuación.
    \item Permitir el análisis manual de correos y la carga de archivos \texttt{.eml}.
    \item Generar reportes PDF y mantener historial local de análisis.
    \item Proporcionar una interfaz web mínima que permita demostrar el funcionamiento.
    \item Evaluar el sistema mediante pruebas y métricas de clasificación.
\end{itemize}

\section{Justificación de las técnicas seleccionadas}

Se seleccionaron cuatro técnicas porque cada una observa el problema desde una perspectiva distinta. La red neuronal permite aprender relaciones no lineales entre características numéricas. Naive Bayes es adecuado para texto porque estima probabilidades a partir de palabras frecuentes en cada clase. El árbol de decisión ofrece una clasificación por umbrales, útil para interpretar patrones. El sistema experto incorpora conocimiento humano de ciberseguridad que no depende únicamente de entrenamiento.

Esta combinación mejora la explicación del resultado final. Si todos los modelos coinciden, la decisión tiene mayor respaldo. Si un modelo detecta riesgo pero otro no, el sistema aún puede mostrar las razones específicas, como fallos de autenticación o presencia de términos sospechosos. Además, la combinación permite cumplir y superar el requisito del enunciado, que solicita al menos tres técnicas de inteligencia artificial.

Otra razón para elegir estas técnicas es que son apropiadas para una implementación manual. Cada algoritmo puede programarse desde sus fundamentos: inicialización de pesos y retropropagación en la red neuronal, conteo de palabras y suavizado de Laplace en Naive Bayes, cálculo de impureza Gini en el árbol, y reglas ponderadas en el sistema experto. Esto permite demostrar dominio del código durante la defensa oral.

\section{Marco teórico}

\subsection{Phishing e ingeniería social}

El phishing es una técnica de fraude digital que busca inducir a una víctima a realizar una acción insegura. Esta acción puede consistir en abrir un enlace, descargar un archivo, introducir credenciales, confirmar datos personales o realizar un pago. En muchos casos, el atacante utiliza mensajes que imitan bancos, universidades, plataformas de correo, empresas de paquetería o servicios de entretenimiento.

Las señales comunes de phishing incluyen lenguaje de urgencia, amenazas de suspensión, solicitudes de contraseña, enlaces a dominios extraños, uso de direcciones IP en lugar de dominios, archivos adjuntos con extensiones ejecutables y errores en mecanismos de autenticación de correo. Aunque ninguna señal individual garantiza que un correo sea malicioso, la acumulación de señales aumenta el riesgo.

\subsection{Clasificación supervisada}

La clasificación supervisada es una tarea de aprendizaje automático en la que se entrena un modelo con ejemplos etiquetados. Cada ejemplo contiene datos de entrada y una etiqueta esperada. En este proyecto, la etiqueta \texttt{1} representa phishing y la etiqueta \texttt{0} representa correo legítimo.

Durante el entrenamiento, los modelos aprenden patrones a partir de los ejemplos. Durante la predicción, reciben un nuevo correo o URL y devuelven una probabilidad o decisión. Para validar el rendimiento, los datos se separan en conjunto de entrenamiento y conjunto de prueba.

\subsection{Red neuronal artificial}

Una red neuronal artificial está formada por unidades llamadas neuronas, organizadas en capas. Cada neurona calcula una suma ponderada de sus entradas y aplica una función de activación. En este proyecto se utiliza la función sigmoide, definida como:

\[
\sigma(x) = \frac{1}{1 + e^{-x}}
\]

La sigmoide produce valores entre 0 y 1, lo que permite interpretar la salida como probabilidad de phishing. La red implementada tiene una capa de entrada, una capa oculta y una neurona de salida. El entrenamiento se realiza mediante retropropagación, ajustando pesos y sesgos para reducir el error entre la predicción y la etiqueta real.

\subsection{Naive Bayes}

Naive Bayes es un clasificador probabilístico basado en el teorema de Bayes. En clasificación textual, se usa para estimar la probabilidad de una clase dado un conjunto de palabras. Su nombre incluye la palabra ``naive'' porque asume independencia condicional entre términos, una simplificación que en la práctica suele funcionar bien para texto.

La idea general es calcular:

\[
P(C \mid X) = \frac{P(X \mid C)P(C)}{P(X)}
\]

Donde \(C\) es la clase y \(X\) representa las palabras del mensaje. Para evitar probabilidades cero cuando aparece una palabra nueva, se utiliza suavizado de Laplace. En el proyecto, Naive Bayes trabaja con tokens extraídos del asunto, cuerpo, URL, remitente y nombres de adjuntos.

\subsection{Árbol de decisión}

Un árbol de decisión clasifica ejemplos mediante divisiones sucesivas de las características. Cada nodo interno evalúa una condición, por ejemplo si una característica es menor que cierto umbral. Cada hoja contiene una decisión o probabilidad.

Para construir el árbol se busca la división que mejor separe las clases. SentinelMail AI utiliza impureza Gini, que mide qué tan mezcladas están las clases en un grupo. Una impureza baja indica que el grupo contiene principalmente ejemplos de una clase.

\[
Gini = 1 - p^2 - (1-p)^2
\]

Donde \(p\) es la proporción de ejemplos positivos. El árbol implementado limita su profundidad para evitar crecimiento excesivo y mantener la clasificación interpretable.

\subsection{Sistema experto}

Un sistema experto representa conocimiento humano mediante reglas. A diferencia de los modelos entrenados, no aprende pesos a partir de datos, sino que evalúa condiciones diseñadas por especialistas. En este proyecto, las reglas se basan en señales comunes de phishing y seguridad de correo.

Ejemplos de reglas son:

\begin{itemize}
    \item Si la URL no usa HTTPS, aumentar el riesgo.
    \item Si el dominio es una dirección IP, aumentar el riesgo.
    \item Si el mensaje solicita contraseña o token, aumentar el riesgo.
    \item Si SPF, DKIM o DMARC fallan, aumentar el riesgo.
    \item Si un adjunto tiene doble extensión, aumentar el riesgo.
\end{itemize}

El sistema experto es valioso porque permite explicar directamente al usuario por qué un mensaje fue considerado riesgoso.

\section{Diseño del sistema}

\subsection{Arquitectura general}

El proyecto se organizó siguiendo una separación por capas inspirada en arquitectura limpia. Esta estructura facilita mantener separada la lógica del dominio, los casos de uso, la infraestructura y la interfaz.

\begin{longtable}{p{4cm}p{10cm}}
\toprule
\textbf{Capa} & \textbf{Responsabilidad} \\
\midrule
Dominio & Define entidades principales, resultados y política de clasificación. \\
Aplicación & Coordina el caso de uso principal de análisis. \\
Infraestructura & Implementa modelos de IA, extracción de características, datasets, parser EML, historial y reportes. \\
Presentación & Expone la interfaz web, endpoints HTTP, archivos estáticos y plantillas HTML. \\
\bottomrule
\end{longtable}

La separación permite explicar claramente el sistema durante la defensa. La interfaz no contiene los algoritmos de IA; solo envía datos al caso de uso. El caso de uso no sabe cómo se dibuja la pantalla; únicamente coordina modelos y devuelve un resultado. Los modelos se ubican en infraestructura porque son detalles técnicos de implementación.

\subsection{Flujo de análisis}

El flujo funcional del sistema es el siguiente:

\begin{enumerate}
    \item El usuario ingresa datos del correo o carga un archivo \texttt{.eml}.
    \item El servidor construye un objeto \texttt{EmailAnalysisRequest}.
    \item El extractor convierte el correo en características numéricas y texto tokenizado.
    \item La red neuronal calcula una probabilidad de phishing.
    \item Naive Bayes calcula una probabilidad textual de phishing.
    \item El árbol de decisión calcula una probabilidad según umbrales aprendidos.
    \item El sistema experto suma reglas activadas y produce razones explicables.
    \item La política de riesgo combina los cuatro puntajes.
    \item El sistema muestra decisión, porcentaje, detalles, reglas, métricas y reporte PDF.
\end{enumerate}

\subsection{Módulos principales}

\begin{longtable}{p{5cm}p{9cm}}
\toprule
\textbf{Archivo} & \textbf{Función} \\
\midrule
\texttt{codigo/app.py} & Punto de entrada de la aplicación web. \\
\texttt{domain/entities.py} & Define entidades como solicitud, adjuntos, indicadores y resultado. \\
\texttt{domain/policies.py} & Combina puntajes y decide nivel de riesgo. \\
\texttt{application/ analyze\_email.py} & Caso de uso principal. \\
\texttt{ai/ feature\_extractor.py} & Extrae características numéricas y tokens. \\
\texttt{ai/ neural\_network.py} & Red neuronal manual. \\
\texttt{ai/ naive\_bayes.py} & Clasificador Naive Bayes manual. \\
\texttt{ai/ decision\_tree.py} & Árbol de decisión manual. \\
\texttt{ai/ expert\_system.py} & Sistema experto de reglas. \\
\texttt{email\_parser.py} & Parser de correos \texttt{.eml} y adjuntos. \\
\texttt{external\_dataset.py} & Carga datasets externos. \\
\texttt{reports.py} & Genera reportes PDF. \\
\texttt{server.py} & Servidor HTTP e interfaz. \\
\bottomrule
\end{longtable}

\subsection{Entidades del dominio}

El dominio define estructuras de datos independientes de la interfaz. La entidad principal de entrada es \texttt{EmailAnalysisRequest}, que contiene asunto, cuerpo, URL, remitente, cabeceras, enlaces, adjuntos y nombre de origen. La salida principal es \texttt{AnalysisResult}, que contiene decisión, nivel, puntaje final, puntajes por técnica, razones, indicadores coincidentes, resumen del correo, características y métricas.

Esta estructura permite que el mismo análisis funcione tanto desde la interfaz web como desde el endpoint JSON \texttt{/api/analyze}.

\section{Implementación}

\subsection{Preparación de datos}

El sistema entrena sus modelos al iniciar la aplicación. Primero carga un dataset simulado realista desde \texttt{SimulatedPhishingDataset}. Este dataset contiene ejemplos legítimos y ejemplos de phishing con asuntos, URLs y cuerpos de mensaje. Además, puede cargar datasets externos desde \texttt{codigo/data/datasets} en formatos CSV, JSON, JSONL o EML.

Cada fila del dataset se convierte en dos representaciones:

\begin{itemize}
    \item Un vector numérico de características para red neuronal y árbol.
    \item Un texto tokenizado para Naive Bayes.
\end{itemize}

Luego se realiza una división de entrenamiento y prueba con una proporción de 75\% para entrenamiento y 25\% para prueba. La semilla fija permite que la separación sea reproducible.

\subsection{Extracción de características}

El extractor genera 23 características. Estas características fueron diseñadas para capturar señales comunes de phishing. Algunas son binarias, como si la URL usa HTTPS o si el dominio contiene punycode. Otras son proporciones normalizadas, como longitud de URL, cantidad de enlaces o cantidad de palabras sospechosas.

\begin{longtable}{p{1cm}p{12cm}}
\toprule
\textbf{No.} & \textbf{Característica} \\
\midrule
1 & La URL usa HTTPS. \\
2 & Longitud de la URL normalizada. \\
3 & Cantidad de puntos en la URL. \\
4 & Cantidad de guiones. \\
5 & Uso de dirección IP como dominio. \\
6 & TLD sospechoso. \\
7 & Presencia de arroba en la URL. \\
8 & Cantidad de dígitos en la URL. \\
9 & Cantidad de enlaces. \\
10 & Palabras de urgencia. \\
11 & Palabras relacionadas con credenciales. \\
12 & Palabras financieras. \\
13 & Uso de acortadores. \\
14 & Presencia de formularios HTML. \\
15 & Extensiones peligrosas mencionadas. \\
16 & Proporción de palabras sospechosas. \\
17 & Cantidad de adjuntos. \\
18 & Diferencia entre \texttt{Reply-To} y \texttt{From}. \\
19 & Diferencia entre \texttt{Return-Path} y \texttt{From}. \\
20 & Fallo SPF, DKIM o DMARC. \\
21 & Cantidad de enlaces declarados. \\
22 & Uso de punycode. \\
23 & Señales riesgosas en adjuntos. \\
\bottomrule
\end{longtable}

\subsection{Red neuronal}

La red neuronal se implementó en la clase \texttt{ManualNeuralNetwork}. Su arquitectura es simple pero suficiente para clasificación binaria: una capa de entrada con 23 características, una capa oculta de 10 neuronas y una neurona de salida.

El método \texttt{\_forward} calcula la propagación hacia adelante. Primero obtiene las activaciones de la capa oculta y luego calcula la salida final. El método \texttt{train} ejecuta 900 épocas de entrenamiento y ajusta los pesos mediante retropropagación.

\begin{lstlisting}[style=pythonstyle, caption={Esquema de predicción de la red neuronal}]
hidden = sigmoid(w1 * x + b1)
output = sigmoid(w2 * hidden + b2)
\end{lstlisting}

El error utilizado es pérdida logarítmica binaria. El gradiente de la salida se calcula como la diferencia entre predicción y etiqueta real. Luego se actualizan los pesos de salida y los pesos de la capa oculta.

\subsection{Naive Bayes}

El clasificador \texttt{ManualNaiveBayes} utiliza una bolsa de palabras. Durante el entrenamiento, cuenta cuántos documentos pertenecen a cada clase y cuántas veces aparece cada palabra en cada clase. También construye el vocabulario total.

Durante la predicción, calcula un puntaje logarítmico para cada clase. Se usa suavizado de Laplace sumando 1 al conteo de cada palabra. Al final, los puntajes se transforman en una probabilidad normalizada de phishing.

\begin{lstlisting}[style=pythonstyle, caption={Idea central del suavizado de Laplace}]
count = word_counts[label][token] + 1
probability = count / (total_words[label] + vocabulary_size)
\end{lstlisting}

Este modelo es especialmente útil para detectar vocabulario asociado a phishing, como ``urgente'', ``contraseña'', ``token'', ``tarjeta'', ``premio'', ``bloqueada'' o ``reembolso''.

\subsection{Árbol de decisión}

El árbol de decisión se implementó en \texttt{ManualDecisionTree}. El entrenamiento construye un árbol binario de manera recursiva. En cada nodo se prueban características y umbrales posibles. La mejor división es aquella que minimiza la impureza Gini ponderada de los grupos izquierdo y derecho.

La predicción recorre el árbol desde la raíz hasta una hoja. Si el valor de la característica es menor que el umbral, se avanza hacia la rama izquierda; en caso contrario, hacia la derecha. La hoja devuelve una probabilidad de phishing basada en la proporción de ejemplos positivos.

El árbol tiene una profundidad máxima de 5 y un tamaño mínimo de nodo de 2, lo cual evita un crecimiento excesivo.

\subsection{Sistema experto}

El sistema experto se implementó en \texttt{SecurityExpertSystem}. A diferencia de los modelos anteriores, no se entrena con datos. Evalúa reglas diseñadas con conocimiento de ciberseguridad y suma puntos de riesgo.

Entre las reglas más importantes están:

\begin{itemize}
    \item URL sin HTTPS: suma 15 puntos.
    \item Dirección IP en el dominio: suma 18 puntos.
    \item TLD sospechoso: suma 10 puntos.
    \item Urgencia en el texto: suma 12 puntos.
    \item Solicitud de credenciales: suma 14 puntos.
    \item Señuelo financiero: suma 10 puntos.
    \item Fallo SPF o DKIM: suma 18 puntos.
    \item Fallo DMARC: suma 20 puntos.
    \item Doble extensión en adjuntos: suma 10 puntos.
    \item Posibles macros en adjuntos: suma 12 puntos.
\end{itemize}

El puntaje del sistema experto se limita a un máximo de 100. Además, cada regla activada produce una razón que se muestra al usuario. Esto aporta interpretabilidad.

\subsection{Análisis de archivos EML y adjuntos}

El parser de correos se implementó con la biblioteca estándar \texttt{email}. Cuando el usuario carga un archivo \texttt{.eml}, el sistema extrae asunto, cuerpo en texto plano, contenido HTML, enlaces, cabeceras y adjuntos.

Para cada adjunto se calcula:

\begin{itemize}
    \item Nombre de archivo.
    \item Tipo de contenido.
    \item Tamaño.
    \item Hash SHA-256.
    \item Extensión.
    \item Presencia de doble extensión.
    \item Sospecha de macros.
    \item Señales de ejecutable de Windows.
\end{itemize}

El hash SHA-256 permite identificar el archivo sin ejecutarlo. El sistema no abre ni ejecuta adjuntos, lo cual es importante desde el punto de vista de seguridad.

\subsection{Política de clasificación final}

La política de riesgo combina los puntajes de las cuatro técnicas. La fórmula utilizada es:

\[
R = 0.36N + 0.24B + 0.18A + 0.22E
\]

Donde:

\begin{itemize}
    \item \(N\) es el puntaje de la red neuronal.
    \item \(B\) es el puntaje de Naive Bayes.
    \item \(A\) es el puntaje del árbol de decisión.
    \item \(E\) es el puntaje del sistema experto.
\end{itemize}

La decisión se determina así:

\begin{longtable}{p{4cm}p{5cm}p{5cm}}
\toprule
\textbf{Puntaje final} & \textbf{Nivel} & \textbf{Decisión} \\
\midrule
0 a 39.99 & BAJO & Probablemente legítimo \\
40 a 69.99 & MEDIO & Sospechoso \\
70 a 100 & ALTO & Probable phishing \\
\bottomrule
\end{longtable}

\section{Interfaz de usuario}

El sistema incluye una interfaz web local ejecutada mediante un servidor HTTP en Python. Para iniciar la aplicación se ejecuta:

\begin{lstlisting}[style=pythonstyle]
python codigo/app.py
\end{lstlisting}

Luego se abre el navegador en:

\begin{center}
\texttt{http://127.0.0.1:8000}
\end{center}

La interfaz permite ingresar asunto, URL, cuerpo del correo, remitente, \texttt{Reply-To}, \texttt{Return-Path} y \texttt{Authentication-Results}. También permite cargar archivos \texttt{.eml}. Al analizar un correo, la aplicación muestra el resultado final, los puntajes por técnica, las razones encontradas, indicadores coincidentes, resumen de cabeceras, adjuntos detectados, métricas del modelo y un enlace al reporte PDF generado.

Además, la aplicación cuenta con rutas de ejemplo:

\begin{itemize}
    \item \texttt{/sample/phishing}
    \item \texttt{/sample/legit}
    \item \texttt{/sample-eml/phishing-completo}
\end{itemize}

También existe un endpoint JSON en:

\begin{center}
\texttt{http://127.0.0.1:8000/api/analyze}
\end{center}

Esto permite demostrar que la lógica del sistema puede ser usada tanto por la interfaz visual como por otros clientes.

\section{Resultados y pruebas}

\subsection{Pruebas automatizadas}

El proyecto incluye pruebas de humo en \texttt{codigo/tests/test\_detector.py}. Estas pruebas validan tres comportamientos importantes:

\begin{itemize}
    \item Un ejemplo de phishing debe clasificarse como riesgo alto.
    \item Un ejemplo legítimo debe clasificarse como riesgo bajo.
    \item Un archivo \texttt{.eml} debe parsearse correctamente, incluyendo adjuntos y autenticación.
\end{itemize}

La ejecución de pruebas produjo el siguiente resultado:

\begin{verbatim}
Ran 3 tests in 1.630s
OK
\end{verbatim}

Esto demuestra que los casos principales del detector funcionan correctamente.

\subsection{Métricas del entrenamiento}

Durante la verificación del sistema se obtuvieron las siguientes métricas:

\begin{longtable}{p{7cm}p{5cm}}
\toprule
\textbf{Métrica} & \textbf{Valor} \\
\midrule
Total de ejemplos & 60 \\
Ejemplos de entrenamiento & 45 \\
Ejemplos de prueba & 15 \\
Ejemplos simulados & 60 \\
Ejemplos externos & 0 \\
Exactitud de entrenamiento & 100.0\% \\
Exactitud de prueba & 100.0\% \\
Verdaderos positivos & 9 \\
Verdaderos negativos & 6 \\
Falsos positivos & 0 \\
Falsos negativos & 0 \\
\bottomrule
\end{longtable}

La exactitud alta se explica porque el dataset educativo contiene señales claras y fue diseñado para demostrar el funcionamiento del sistema. En un entorno real, se esperaría que la exactitud varíe y que sea necesario entrenar con corpus más amplios y diversos.

\subsection{Caso de prueba: phishing}

Se probó el siguiente caso:

\begin{itemize}
    \item Asunto: Cuenta bloqueada urgente.
    \item URL: \texttt{http://seguridad-banco.example.verify-login.ru/acceso}
    \item Mensaje: Urgente: su cuenta será suspendida. Verifique usuario, contraseña y token.
\end{itemize}

El resultado fue:

\begin{longtable}{p{7cm}p{5cm}}
\toprule
\textbf{Campo} & \textbf{Valor} \\
\midrule
Decisión & Probable phishing \\
Nivel & ALTO \\
Riesgo final & 98.02\% \\
Red neuronal & 100.0\% \\
Naive Bayes & 100.0\% \\
Árbol de decisión & 100.0\% \\
Sistema experto & 91.0\% \\
\bottomrule
\end{longtable}

Las principales razones activadas fueron: la URL no usa HTTPS, el dominio usa terminación sospechosa, el dominio contiene subdominios o guiones, el mensaje usa urgencia y solicita credenciales.

\subsection{Caso de prueba: correo legítimo}

Se probó el siguiente caso:

\begin{itemize}
    \item Asunto: Aviso de mantenimiento.
    \item URL: \texttt{https://hosting.example/status}
    \item Mensaje: El servicio tendrá una ventana de mantenimiento programada el sábado.
\end{itemize}

El resultado fue:

\begin{longtable}{p{7cm}p{5cm}}
\toprule
\textbf{Campo} & \textbf{Valor} \\
\midrule
Decisión & Probablemente legítimo \\
Nivel & BAJO \\
Riesgo final & 0.0\% \\
Red neuronal & 0.01\% \\
Naive Bayes & 0.0\% \\
Árbol de decisión & 0.0\% \\
Sistema experto & 0.0\% \\
\bottomrule
\end{longtable}

El sistema no activó reglas críticas de phishing, lo cual coincide con el comportamiento esperado.

\section{Cumplimiento del enunciado}

El enunciado solicita una aplicación funcional que resuelva un problema específico utilizando al menos tres técnicas de inteligencia artificial cubiertas en el curso. SentinelMail AI cumple este requisito porque resuelve el problema de detección de phishing y aplica cuatro técnicas diferentes.

\begin{longtable}{p{7cm}p{7cm}}
\toprule
\textbf{Requisito del enunciado} & \textbf{Cumplimiento en el proyecto} \\
\midrule
Resolver un problema claramente definido con IA & Se resuelve la detección de phishing en correos y URLs. \\
Implementar al menos 3 técnicas de IA & Se implementan 4: red neuronal, Naive Bayes, árbol de decisión y sistema experto. \\
Funcionar con datos reales o simulados realistas & Usa dataset simulado realista y puede cargar datasets externos. \\
Incluir interfaz de usuario mínima & Incluye interfaz web local. \\
Algoritmos centrales implementados manualmente & Los modelos se implementaron sin bibliotecas de aprendizaje automático. \\
Documentación técnica clara & Este informe documenta problema, diseño, implementación, resultados y uso. \\
Resultados verificados & Se ejecutaron pruebas automatizadas y casos manuales. \\
Manual de usuario & Se incluye en una sección específica. \\
\bottomrule
\end{longtable}

En cuanto al formato de entrega, el enunciado solicita una carpeta de código, una carpeta de documentación, una carpeta de presentación y un archivo README. El repositorio contiene la carpeta \texttt{codigo}, la carpeta \texttt{documentacion} y el archivo \texttt{README.txt}. Para la entrega final, se recomienda agregar una carpeta \texttt{presentacion} con las diapositivas de la defensa y comprimir todo en formato ZIP con el nombre solicitado por el profesor.

\section{Limitaciones}

Aunque el sistema cumple los objetivos académicos, tiene limitaciones importantes:

\begin{itemize}
    \item El dataset incluido es simulado, aunque realista. Esto facilita la demostración, pero no sustituye un corpus real amplio.
    \item No consulta servicios externos de reputación de dominios.
    \item No ejecuta adjuntos en sandbox; solo realiza análisis estático.
    \item Puede generar falsos positivos si un correo legítimo usa lenguaje urgente.
    \item Puede generar falsos negativos si un phishing está escrito de forma muy cuidadosa y sin señales evidentes.
    \item El reporte PDF es funcional, pero básico en diseño visual.
\end{itemize}

Estas limitaciones no invalidan el proyecto; más bien delimitan su alcance académico. El objetivo no es reemplazar una plataforma profesional de seguridad, sino demostrar la aplicación de técnicas de inteligencia artificial a un problema realista.

\section{Mejoras futuras}

Como mejoras futuras se proponen:

\begin{itemize}
    \item Entrenar con datasets públicos más grandes y diversos.
    \item Agregar validación con dominios reales y reputación en línea.
    \item Incorporar análisis de similitud visual para detectar dominios homógrafos.
    \item Analizar adjuntos en un entorno sandbox controlado.
    \item Mejorar el diseño del reporte PDF.
    \item Agregar explicabilidad por característica para la red neuronal.
    \item Permitir exportar resultados en CSV o JSON desde la interfaz.
    \item Agregar autenticación de usuarios si se desea usar el sistema en una red institucional.
\end{itemize}

\section{Conclusiones}

SentinelMail AI demuestra que es posible construir una aplicación funcional de detección de phishing utilizando técnicas de inteligencia artificial implementadas manualmente. El proyecto integra cuatro enfoques diferentes y los combina en una política de riesgo clara. Esta integración permite analizar texto, características numéricas, reglas de seguridad y estructura de correos electrónicos.

La red neuronal aprende patrones numéricos asociados al riesgo. Naive Bayes identifica vocabulario frecuente en mensajes de phishing. El árbol de decisión aporta clasificación por umbrales e interpretación estructurada. El sistema experto incorpora reglas explícitas de ciberseguridad y genera explicaciones comprensibles para el usuario.

El sistema cumple con los requisitos técnicos principales del enunciado: tiene un problema definido, implementa más de tres técnicas de IA, funciona con datos simulados realistas, posee interfaz web, incluye pruebas, genera reportes y mantiene historial de análisis. Además, los algoritmos centrales fueron programados sin bibliotecas externas de aprendizaje automático, lo cual evidencia comprensión de los fundamentos.

Como aprendizaje principal, el proyecto muestra que la inteligencia artificial aplicada no depende únicamente de modelos complejos. En problemas de seguridad, la combinación de modelos estadísticos, aprendizaje supervisado y reglas expertas puede ofrecer resultados más explicables y útiles que una única técnica aislada.

\section{Manual de usuario}

\subsection{Requisitos}

Para ejecutar el sistema se requiere:

\begin{itemize}
    \item Python 3.10 o superior.
    \item Sistema operativo Windows, Linux o macOS.
    \item Navegador web moderno.
    \item No se requieren bibliotecas externas para ejecutar la aplicación.
\end{itemize}

\subsection{Instalación}

El usuario debe descomprimir el proyecto y abrir una terminal en la carpeta raíz. La estructura esperada incluye:

\begin{itemize}
    \item \texttt{codigo/}
    \item \texttt{documentacion/}
    \item \texttt{README.txt}
\end{itemize}

\subsection{Ejecución}

Desde la carpeta raíz del proyecto, ejecutar:

\begin{verbatim}
python codigo/app.py
\end{verbatim}

El sistema mostrará un mensaje similar a:

\begin{verbatim}
Detector iniciado en http://127.0.0.1:8000
Presione Ctrl+C para detener.
\end{verbatim}

Luego se debe abrir el navegador en:

\begin{verbatim}
http://127.0.0.1:8000
\end{verbatim}

\subsection{Uso básico}

Para analizar un correo manualmente:

\begin{enumerate}
    \item Escribir o pegar el asunto.
    \item Ingresar la URL principal.
    \item Pegar el cuerpo del mensaje.
    \item Opcionalmente completar remitente, \texttt{Reply-To}, \texttt{Return-Path} y resultados de autenticación.
    \item Presionar el botón de análisis.
    \item Revisar el nivel de riesgo, puntajes por técnica y razones encontradas.
\end{enumerate}

\subsection{Carga de archivos EML}

Para analizar un correo real exportado:

\begin{enumerate}
    \item Seleccionar un archivo con extensión \texttt{.eml}.
    \item Enviar el formulario.
    \item El sistema extraerá asunto, cuerpo, enlaces, cabeceras y adjuntos.
    \item Se mostrará el resultado de riesgo.
\end{enumerate}

\subsection{Reportes e historial}

Cada análisis genera un reporte PDF en la carpeta \texttt{codigo/data/reports}. También se guarda un resumen en \texttt{codigo/data/history/analysis\_history.json}. El historial conserva los últimos análisis para consulta desde la interfaz.

\subsection{Pruebas}

Para ejecutar las pruebas automatizadas:

\begin{verbatim}
python -m unittest discover codigo/tests
\end{verbatim}

El resultado esperado es que las pruebas finalicen con \texttt{OK}.

\section{Guía para la defensa oral}

Para una exposición de 15 a 20 minutos, se recomienda dividir la presentación de la siguiente manera:

\begin{longtable}{p{4cm}p{10cm}}
\toprule
\textbf{Tiempo aproximado} & \textbf{Contenido} \\
\midrule
2 minutos & Presentación del problema de phishing y motivación. \\
2 minutos & Objetivo general y alcance del sistema. \\
4 minutos & Explicación de las cuatro técnicas de IA. \\
3 minutos & Arquitectura y flujo de análisis. \\
4 minutos & Demostración en vivo con caso phishing y caso legítimo. \\
2 minutos & Resultados, métricas y pruebas. \\
2 minutos & Conclusiones, limitaciones y mejoras futuras. \\
\bottomrule
\end{longtable}

Una respuesta breve para explicar el proyecto podría ser:

\begin{quote}
SentinelMail AI es una aplicación web local que analiza correos electrónicos y URLs para estimar riesgo de phishing. El sistema combina cuatro técnicas de inteligencia artificial implementadas manualmente: red neuronal, Naive Bayes, árbol de decisión y sistema experto. Cada técnica evalúa señales distintas y luego una política de riesgo combina los puntajes para producir una decisión final explicable.
\end{quote}

Si el profesor pregunta por qué se usan cuatro técnicas, la respuesta sugerida es:

\begin{quote}
Porque el phishing no se detecta con una sola señal. La red neuronal aprende patrones numéricos, Naive Bayes analiza lenguaje, el árbol clasifica por umbrales y el sistema experto incorpora reglas de ciberseguridad. Al combinarlas, el sistema obtiene una evaluación más completa y explicable.
\end{quote}

Si el profesor pregunta cómo se programaron los algoritmos, la respuesta sugerida es:

\begin{quote}
Los algoritmos centrales se implementaron manualmente. La red neuronal usa sigmoide, pérdida logarítmica y retropropagación; Naive Bayes usa conteo de tokens, probabilidad por clase y suavizado de Laplace; el árbol usa divisiones por umbral e impureza Gini; el sistema experto usa reglas ponderadas de seguridad.
\end{quote}

\section{Referencias bibliográficas}

\begin{thebibliography}{9}

\bibitem{russell}
Russell, S. y Norvig, P. (2021). \textit{Artificial Intelligence: A Modern Approach}. Pearson.

\bibitem{mitchell}
Mitchell, T. M. (1997). \textit{Machine Learning}. McGraw-Hill.

\bibitem{bishop}
Bishop, C. M. (2006). \textit{Pattern Recognition and Machine Learning}. Springer.

\bibitem{goodfellow}
Goodfellow, I., Bengio, Y. y Courville, A. (2016). \textit{Deep Learning}. MIT Press.

\bibitem{owasp}
OWASP Foundation. (2024). \textit{Phishing and social engineering security awareness resources}. OWASP.

\bibitem{nist}
National Institute of Standards and Technology. (2020). \textit{Security and Privacy Controls for Information Systems and Organizations}. NIST Special Publication 800-53.

\bibitem{rfc}
Kitterman, S. (2014). \textit{Sender Policy Framework (SPF) for Authorizing Use of Domains in Email}. RFC 7208.

\end{thebibliography}

\end{document}
